\documentclass[10pt,twocolumn,twoside,a4paper]{report}

\usepackage[slovak,english]{babel}
\usepackage[IL2]{fontenc}
\usepackage[utf8]{inputenc}
\usepackage{graphicx}
\usepackage{url}
\usepackage{hyperref}

\usepackage{cite}

\graphicspath{ {./images/} }

\pagestyle{headings}

\title{Personalized Recommendation Systems in Social Media\thanks{Semester project for the course Engineering Methods, academic year 2024/25, supervised by Ing. Richard Marko, PhD.}} % name and title of the course supervisor

\author{Yevhenii Deshchenko\\[2pt]
	{\small Slovak Technical University in Bratislava}\\
	{\small Faculty of Informatics and Information Technology}\\
	{\small \texttt{xdeshchenko@stuba.sk}}
	}

\date{\small 10 October 2024}

\begin{document}

\maketitle

\begin{abstract}
With the growing popularity of social media, recommendation systems have become increasingly complex and sophisticated. 
These systems rely on a variety of algorithms, including those based on graph theory, statistics, machine and deep learning. 
Content-based recommendation systems vary significantly across many different social media platforms, as they depend on the 
type of primary content users share and the options provided for user interaction with the content. 
Most approaches use as much data as available, leading to massive data collection that raises multiple concerns, such as privacy issues.


As social media platforms continue to grow in popularity, the inner workings of these algorithms are often kept secret,
as they have become valuable tools for third parties seeking to exploit them for visibility and influence.
This exploitation frequently aims to push content with political or other agendas, in an effort to gain popularity on these platforms.
However, there are exceptions, such as Twitter, which made its recommendation algorithm public after a data breach.
This decision led to significant criticism from its user base. The secrecy and potential for manipulation continue to
spark debates about the role and transparency of recommendation systems in shaping user experiences and content visibility on social media.
\end{abstract}

% \includegraphics[width=\textwidth]{fiit.eps}

\section{Introduction}

In a modern era where social media has become an integral part of daily life, personalized recommendation systems play a
crucial role in enhancing marketing, entertainment, and the overall user experience. By predicting and delivering content
that resonates with users, these systems contribute significantly to social media growth, fostering higher user engagement
and supporting an interconnected ecosystem of content creators, businesses, and users alike.

While these systems are built upon statistical and mathematical foundations, practical implementations vary widely.
They often combine several approaches, including complex algorithms, databases, big data analytics, machine learning,
and deep learning techniques, to optimize predictions and deliver tailored content. 

However, the reliance on extensive user data raises privacy concerns, as large amounts of personal information must
be stored and processed. Additionally, biases within these algorithms can lead to unintended consequences, potentially
skewing recommendations or enabling exploitation. This paper will explore the methodologies behind recommendation systems,
their societal impacts, and the balance between personalization and ethical considerations.

\section{Some Section} \label{some_section}

\section{Conclusion} \label{conclusion}

%\acknowledgement{If you want to thank someone\ldots}

\nocite{*}
\bibliography{literatura}
\bibliographystyle{plain} % or alpha, abbrv, or any other style
\end{document}
